%!TEX root = ../template.tex
%%%%%%%%%%%%%%%%%%%%%%%%%%%%%%%%%%%%%%%%%%%%%%%%%%%%%%%%%%%%%%%%%%%%
%% chapter6.tex
%% NOVA thesis document file
%%
%% Chapter with lots of dummy text
%%%%%%%%%%%%%%%%%%%%%%%%%%%%%%%%%%%%%%%%%%%%%%%%%%%%%%%%%%%%%%%%%%%%

\typeout{NT FILE chapter6.tex}%

\chapter{Planning}
\label{cha:planning}


\section{Introduction}

In this chapter, the work plan for this dissertation is presented, considering a period of 24 weeks (6 months).

The objective of this work is to add support for attribute grammars to the system OCamlFLAT/OFLAT. The implementation work aims to extend existing software, which already includes some relevant elements, such as support for context-free grammars and basic visualization and graphical interaction mechanisms.

\section{Support for Attribute Grammars in the OCamlFLAT Library (8 weeks)}

At the library level, several functionalities are planned:

\begin{itemize}
    \item Representation of attributes and their calculation rules within grammars.
    \item Circularity testing of synthesized and inherited attributes in a grammar.
    \item Calculation of synthesized and inherited attributes on an already constructed derivation tree.
    \item Calculation of synthesized and inherited attributes during the derivation tree construction of a word, i.e., during parsing.
    \item Introduction of "conditions" in grammars to use attribute values to restrict the generated language.
    \item Integration of attribute grammars into the existing support for composite models in the library. This involves attempting various operations on attribute grammars: union, intersection, difference, concatenation, and Kleene closure.
    \item Integration of grammars into the existing support for pedagogical exercises in the library.
\end{itemize}

\section{Support for Grammars in the OFLAT Web Application (6 weeks)}

Attribute grammars must be integrated into the existing graphical interface. The additions should have pedagogical features, and some interesting ideas may emerge for the graphical interface.

There is already support for context-free grammars in the graphical application. That part of the code will be duplicated, refactored, and adapted. Understanding and modifying the graphical application code is not a trivial task.

What should be supported in the graphical interface for attribute grammars:

\begin{itemize}
    \item Visualization.
    \item Editing.
    \item Visualization of the results of operations available in the library.
    \item Visualization of derivation trees annotated with attribute calculations. In a more ambitious version, attribute calculation could be animated.
    \item Integration into the existing support for composite models in the graphical application.
\end{itemize}

\section{Pedagogical Exercises (1 week)}

Create a collection of interesting pedagogical examples and encode them into a set of exercises in the format supported by the OCamlFLAT/OFLAT system.

\section{Evaluation (1 week)}

\begin{itemize}
    \item The correctness of the newly implemented code will be assessed through a set of unit tests, which will be written alongside code development.
    \item Analyzing the system's limitations by creating examples that pose challenges due to size or execution time.
    \item Seeking feedback from potential users of the system, particularly regarding the graphical interface. For example, teachers and students of the TC course.
\end{itemize}

\section{Dissertation Writing (8 weeks)}

Ideally conducted in small increments, concurrently with other tasks. Two months are considered for the production of the final document. It is important to have a near-final version ready two weeks before the submission deadline to facilitate inviting an examiner.

\section{Gantt Chart (Table Representation)}

\begin{table}[h]
    \centering
    \begin{tabular}{|l|c|}
        \hline
        \textbf{Task} & \textbf{Weeks} \\
        \hline
        OCamlFLAT Library Support & 1 - 8 \\
        OFLAT Web Application Support & 9 - 14 \\
        Pedagogical Exercises & 15 \\
        Evaluation & 16 \\
        Dissertation Writing & 17 - 24 \\
        \hline
    \end{tabular}
    \caption{Work Plan Timeline}
\end{table}